\begin{theorem}[Global vanishing from one degree-two vanishing]
\hypertarget{tag-03J0}{}\label{wbII:thm:corrected-cdp-global}
\StatementRecord{wbII:thm:corrected-cdp-global}{proved}{03J0}{\SeriesRef{wbII:lem:cdp-generic-endpoints}, \SeriesRef{wbII:eq:W:5.28}}
\label{wbII:thm:corrected-global}
Let \(X\) be a smooth compact connected complex threefold such that
\[
 H^1(X;\mathbb Z)=H^2(X;\mathbb Z)=0,\qquad a(X)=1,
\]
and suppose its algebraic reduction is a holomorphic map
\(f:X\to C\).  If there exists \(L_0\in\operatorname{Pic}(X)\) with
\[
 H^2(X,T_X\otimes L_0)=0,
\tag{W.51}\hypertarget{tag-03HM}{}\label{wbII:eq:W:5.24}
\]
then, for \(M\) in a dense open subset of \(\operatorname{Pic}(X)\),
\[
 H^i(X,T_X\otimes M)=0\quad(i\ne1),
\tag{W.52}\hypertarget{tag-03HN}{}\label{wbII:eq:W:5.25}
\]
\[
 \chi(X,T_X\otimes M)\le0,\qquad
 \left\langle c_3(T_X),[X]\right\rangle\le0.
\tag{W.53}\hypertarget{tag-03HO}{}\label{wbII:eq:W:5.26}
\]
\end{theorem}



\begin{proof}

\begin{enumerate}[label=\arabic*.,leftmargin=*]
\item The exponential sequence and the two integral vanishings give
\[
 \operatorname{Pic}(X)=\operatorname{Pic}^0(X)
 \simeq H^1(X,\mathcal O_X).
\tag{W.54}\hypertarget{tag-03HP}{}\label{wbII:eq:W:5.27}
\]
Thus this Picard group is a connected irreducible complex vector group.
Choose additive \v{C}ech cocycles $a^{(r)}_{ij}$ representing a basis of
$H^1(X,\mathcal O_X)$. For coordinates $\lambda_r$ on that vector space,
\[
 g_{ij}(\lambda)=\exp\!\left(\sum_r\lambda_r a^{(r)}_{ij}\right)
\]
satisfy the multiplicative cocycle condition and define a holomorphic
line bundle on $X\times H^1(X,\mathcal O_X)$. Its restriction at every
parameter represents the corresponding line bundle, and the projection to the parameter space is proper.  Any fibre divisor \(D\) of
\(f\) is nonzero and effective, and its line bundle lies in
\(\operatorname{Pic}^0(X)\) by \(H^2(X;\mathbb Z)=0\).  Applying
\cref{wbII:lem:cdp-generic-endpoints} to \(B=T_X\) shows that the degree-zero and
degree-three groups in \eqref{wbII:eq:W:5.25} vanish away from two proper analytic subsets
of \(\operatorname{Pic}(X)\).

\item The cohomology-jump locus
\[
 \{M\in\operatorname{Pic}(X):h^2(X,T_X\otimes M)\ge1\}
\]
is locally determinantal in this family
\cite[Chapter~III, \S4, Theorem~4.1, p.~124]{BS76push}.  It is proper because \eqref{wbII:eq:W:5.24} gives a
point outside it.  Its complement meets the two endpoint
vanishing loci in a dense open set, where \eqref{wbII:eq:W:5.25} holds and
\[
 \chi(X,T_X\otimes M)=-h^1(X,T_X\otimes M)\le0.
\tag{W.55}\hypertarget{tag-03HQ}{}\label{wbII:eq:W:5.28}
\]

\item Every \(M\in\operatorname{Pic}(X)\) has rationally trivial first Chern
class.  The hypotheses also give \(c_1(T_X)=0\) in rational cohomology and,
by Poincar\'e duality, \(H^4(X;\mathbb Q)=0\), so
\(c_2(T_X)=0\) rationally.  The degree-six term of
\(\operatorname{ch}(T_X\otimes M)\operatorname{td}(X)\) is consequently
\(\tfrac12c_3(T_X)\).  Hirzebruch--Riemann--Roch and \eqref{wbII:eq:W:5.28} give the inequalities
in \eqref{wbII:eq:W:5.26}:
\[
 \chi(X,T_X\otimes M)
 =\frac12\left\langle c_3(T_X),[X]\right\rangle\le0.
\]

\end{enumerate}
\end{proof}